% Discrete Mathematics (MAL2190) - IIT Jodhpur
% Project Viva-Voce Defense & Comprehensive Mathematical Reasoning Guide
% Course: Discrete Mathematics (MAL2190)
% Group 14: Dungar Soni, Radtiya Saraf, Prakhar Goyal, Arafat, Bhawani

\documentclass[11pt,a4paper]{article}
\usepackage[T1]{fontenc}
\usepackage{lmodern}
\usepackage[a4paper,margin=0.9in,headheight=15pt]{geometry}
\usepackage{amsmath,amssymb,amsthm}
\usepackage{graphicx}
\usepackage[table]{xcolor}
\usepackage{array,tabularx,longtable,booktabs}
\usepackage{enumitem}
\usepackage{fancyhdr,lastpage}
\usepackage{titlesec}
\usepackage[hidelinks]{hyperref}
\usepackage{mdframed}

% Custom Styling & Palette
\definecolor{primaryblue}{RGB}{24, 76, 120}
\definecolor{accentteal}{RGB}{0, 128, 128}
\definecolor{darkcharcoal}{RGB}{33, 37, 41}
\definecolor{boxbg}{RGB}{245, 248, 252}
\definecolor{boxborder}{RGB}{180, 205, 230}
\definecolor{highlightbg}{RGB}{254, 249, 231}
\definecolor{highlightborder}{RGB}{241, 196, 15}

\setlength{\parindent}{0pt}
\setlength{\parskip}{6pt}
\renewcommand{\arraystretch}{1.3}
\setlist{leftmargin=*,itemsep=3pt,topsep=3pt,parsep=0pt}

\titleformat{\section}{\large\bfseries\color{primaryblue}}{\thesection.}{0.5em}{}
\titleformat{\subsection}{\normalsize\bfseries\color{accentteal}}{\thesubsection}{0.5em}{}
\titleformat{\subsubsection}{\small\bfseries\color{darkcharcoal}}{\thesubsubsection}{0.5em}{}

\pagestyle{fancy}
\fancyhf{}
\fancyhead[L]{\small\color{primaryblue}\textbf{MAL2190: Discrete Mathematics} --- Viva Defense Guide}
\fancyhead[R]{\small\color{primaryblue}\textbf{Group 14} (IIT Jodhpur)}
\fancyfoot[L]{\footnotesize Indian Institute of Technology Jodhpur}
\fancyfoot[R]{\footnotesize Page \thepage\ of \pageref{LastPage}}
\renewcommand{\headrulewidth}{0.5pt}
\renewcommand{\footrulewidth}{0.5pt}

\newmdenv[
  backgroundcolor=boxbg,
  linecolor=boxborder,
  linewidth=1.2pt,
  roundcorner=5pt,
  innertopmargin=8pt,
  innerbottommargin=8pt,
  innerleftmargin=10pt,
  innerrightmargin=10pt,
  skipabove=8pt,
  skipbelow=8pt
]{conceptbox}

\newmdenv[
  backgroundcolor=highlightbg,
  linecolor=highlightborder,
  linewidth=1.2pt,
  roundcorner=5pt,
  innertopmargin=8pt,
  innerbottommargin=8pt,
  innerleftmargin=10pt,
  innerrightmargin=10pt,
  skipabove=8pt,
  skipbelow=8pt
]{vivabox}

\theoremstyle{definition}
\newtheorem{theorem}{Theorem}
\newtheorem{lemma}{Lemma}
\newtheorem{definition}{Definition}

\begin{document}
\thispagestyle{plain}

\begin{center}
  {\Large\bfseries\color{primaryblue} Indian Institute of Technology Jodhpur\par}
  \vspace{3pt}
  {\large\bfseries Department of Computer Science \& Engineering\par}
  \vspace{3pt}
  {\large\bfseries MAL2190: Discrete Mathematics --- Project Viva-Voce Defense Guide\par}
  \vspace{8pt}
  \rule{\textwidth}{1pt}\par
  \vspace{6pt}
  {\LARGE\bfseries Logical Rationale, Mathematical Foundations, and Oral Examination Defense Strategy\par}
  \vspace{4pt}
  {\large\textit{Dual-Realm Discrete Adventures: Monster Sort \& The Robot Maze Challenge}\par}
  \vspace{6pt}
  \rule{\textwidth}{1pt}\par
  \vspace{10pt}
\end{center}

\begin{tabularx}{\textwidth}{@{}P{3.5cm}Y P{3.5cm}Y@{}}
\textbf{Group Number:} & Group 14 & \textbf{Course Code:} & MAL2190 (Discrete Mathematics) \\
\textbf{Academic Year:} & 2026--2027 & \textbf{Faculty Examiner:} & Course Instructor \\
\end{tabularx}

\vspace{6pt}
\textbf{Group Members:}
\begin{enumerate}[label=\textbf{\arabic*.}]
  \item \textbf{Dungar Soni} (B23CS1105) --- Lead Full-Stack Software Architecture \& Systems Integration
  \item \textbf{Radtiya Saraf} (B23CS1107) --- Theoretical Discrete Mathematics \& Formal Mathematical Proofs
  \item \textbf{Prakhar Goyal} (B23CS1106) --- Planar Geometric Modeling, SVG Topology \& Grid Graph Physics
  \item \textbf{Arafat} (B23CS1103) --- Procedural Feature Schemas, Vector Visuals \& Audio Synthesis
  \item \textbf{Bhawani} (B23CS1104) --- Pedagogical Psychology, Personality Quiz Logic \& Maze Experimentation
\end{enumerate}

\hrulefill

\section{Logical Rationale: Why This Project? (Selection Logic)}

When the faculty examiner asks: \textit{``What was the logic behind selecting this project, and why did you design it this way?''}, the answer must articulate both **pedagogical necessity** and **mathematical depth**.

\subsection{The Pedagogical Problem with Traditional Discrete Mathematics Teaching}
Discrete Mathematics is traditionally taught as a collection of disjoint formalisms:
\begin{enumerate}
  \item \textbf{Abstract Detachment:} Concepts such as characteristic functions $\chi_A(x)$, Venn partitions, minterm expansions, De Morgan's laws, and graph search invariants are presented via static formulas and truth tables on chalkboards. Students struggle to develop spatial and operational intuition.
  \item \textbf{Compartmentalization:} Set Theory, Boolean Algebra, and Graph Theory are often taught in isolated textbook chapters. Students rarely realize that \textbf{graph traversal is dynamically equivalent to maintaining partitions of a state space}.
  \item \textbf{Passive Reception vs.\ Constructivist Discovery:} Under Seymour Papert's \textit{constructionist} educational theory, learners grasp abstract mathematics most effectively when they actively manipulate objects, hypothesize, encounter constraints, and receive immediate deterministic feedback.
\end{enumerate}

\subsection{Why a Dual-Realm Ecosystem? (Monster Sort + Robot Maze)}
Rather than building an isolated, trivial puzzle, we created a unified, complementary dual-realm system:
\begin{itemize}
  \item \textbf{Monster Sort (The Static, Taxonomic Realm):} Focuses on \textbf{classification, unary predicates, Boolean set algebra, and Venn partitioning}. It answers: \textit{``Given static objects with discrete features, into which unique spatial block of the partition do they belong?''}
  \item \textbf{Robot Maze (The Dynamic, Algorithmic Realm):} Focuses on \textbf{relations, grid graphs, state-space navigation, and geodesic shortest paths via BFS}. It answers: \textit{``Given a state-space graph, what sequence of discrete relations transitions an agent from source to target with minimal cost?''}
  \item \textbf{The Sorting Hat Gateway:} Connects the two domains through cognitive profiling. An analytical classifier naturally excels at Boolean feature extraction (Monster Sort), whereas a sequential explorer naturally excels at graph pathfinding (Robot Maze). Profiling students highlights that both paradigms are essential branches of discrete mathematics.
\end{itemize}

\begin{conceptbox}
\textbf{The Core Viva Thesis:}
``Our project demonstrates the fundamental duality of discrete mathematics: \textbf{Monster Sort} explores the \textit{static partition of a universe} using Boolean characteristic predicates, while \textbf{Robot Maze} explores the \textit{dynamic traversal of a graph relation} using Breadth-First Search. Together, they prove that set theory and graph theory are two sides of the same mathematical coin.''
\end{conceptbox}

\newpage
\section{Mathematical Reasoning: Realm 1 --- Monster Sort}

\subsection{Formal Modeling of Monsters and Predicates}
\begin{definition}[Universe of Discourse]
The monster kingdom is modeled as a finite universal set:
$$U = \{M_1, M_2, \dots, M_n\}$$
Each creature $M_j \in U$ is characterized by a 7-dimensional attribute vector:
$$\vec{f}(M_j) = (\text{color}_j, \text{horns}_j, \text{wings}_j, \text{eyes}_j, \text{shape}_j, \text{pattern}_j, \text{size}_j) \in \mathcal{C} \times \mathcal{H} \times \mathcal{W} \times \mathcal{E} \times \mathcal{S} \times \mathcal{P} \times \mathcal{Z}$$
\end{definition}

\begin{definition}[Unary Boolean Predicates \& Truth Sets]
Given $k$ properties of interest, we define $k$ unary predicates $P_i : U \to \{0, 1\}$ for $i \in \{1, \dots, k\}$. Each predicate induces a truth set (subset of $U$):
$$A_i = \{x \in U \mid P_i(x) = 1\} \subseteq U$$
The complement $A_i^c = U \setminus A_i = \{x \in U \mid P_i(x) = 0\}$ contains all creatures lacking property $i$.
\end{definition}

\subsection{Characteristic Indicator Functions \& Boolean Isomorphism}
The mathematical foundation of Monster Sort relies on the **Characteristic Function**:
$$\chi_{A_i}(x) = \begin{cases} 1, & \text{if } x \in A_i \\ 0, & \text{if } x \notin A_i \end{cases}$$
This establishes a rigorous algebraic isomorphism between the power set Boolean algebra $(\mathcal{P}(U), \cup, \cap, {}^c, \emptyset, U)$ and the bit-vector algebra $(\{0, 1\}^U, \lor, \land, \neg, \mathbf{0}, \mathbf{1})$:
\begin{align*}
\chi_{A \cap B}(x) &= \chi_A(x) \land \chi_B(x) = \chi_A(x) \cdot \chi_B(x) \\
\chi_{A \cup B}(x) &= \chi_A(x) \lor \chi_B(x) = \chi_A(x) + \chi_B(x) - \chi_A(x)\chi_B(x) \\
\chi_{A^c}(x) &= \neg \chi_A(x) = 1 - \chi_A(x) \\
\chi_{A \setminus B}(x) &= \chi_{A \cap B^c}(x) = \chi_A(x)(1 - \chi_B(x))
\end{align*}

\subsection{Venn Diagram Topology and Minterm Partitions}
\begin{definition}[Minterm / Atom of Boolean Algebra]
For $k$ sets $A_1, \dots, A_k$, each binary vector $\vec{b} = (b_1, \dots, b_k) \in \{0, 1\}^k$ uniquely identifies a fundamental minterm set:
$$M_{\vec{b}} = \bigcap_{i=1}^k A_i^{(b_i)} \quad \text{where} \quad A_i^{(b_i)} = \begin{cases} A_i, & \text{if } b_i = 1 \\ A_i^c, & \text{if } b_i = 0 \end{cases}$$
\end{definition}

\begin{theorem}[Venn Partition Theorem]
The family $\mathcal{M} = \{M_{\vec{b}} \mid \vec{b} \in \{0, 1\}^k\}$ forms a strict mathematical partition of $U$:
\begin{enumerate}
  \item \textbf{Pairwise Disjointness:} For all $\vec{b} \neq \vec{c}$, $M_{\vec{b}} \cap M_{\vec{c}} = \emptyset$.
  \item \textbf{Exhaustive Covering:} $\bigcup_{\vec{b} \in \{0, 1\}^k} M_{\vec{b}} = U$.
  \item \textbf{Unique Assignment:} Every element $x \in U$ belongs to exactly one minterm $M_{\vec{b}(x)}$ where $\vec{b}(x) = (\chi_{A_1}(x), \dots, \chi_{A_k}(x))$.
\end{enumerate}
\end{theorem}

\begin{proof}[Proof Sketch for Examiner]
\textbf{1. Disjointness:} If $\vec{b} \neq \vec{c}$, they differ at some coordinate $j$ ($b_j = 1, c_j = 0$). Then $M_{\vec{b}} \subseteq A_j$ and $M_{\vec{c}} \subseteq A_j^c$. Since $A_j \cap A_j^c = \emptyset$, their intersection is empty: $M_{\vec{b}} \cap M_{\vec{c}} = \emptyset$.

\textbf{2. Covering:} For any $x \in U$, by the Law of Excluded Middle, either $x \in A_i$ or $x \in A_i^c$ for each $i$. Setting $b_i = \chi_{A_i}(x)$ yields an index $\vec{b}^*$ such that $x \in M_{\vec{b}^*}$. Since $M_{\vec{b}^*} \subseteq \bigcup M_{\vec{b}}$, $U = \bigcup M_{\vec{b}}$.

\textbf{3. Uniqueness:} If $x$ belonged to two different minterms, their intersection would contain $x$, contradicting pairwise disjointness.
\end{proof}

\subsection{Why the Universal Canvas is Mathematically Mandatory}
In traditional classroom Venn diagrams, students frequently ignore the space outside the circles. Mathematically:
$$\mathcal{R}_0 = M_{(0,0,\dots,0)} = \bigcap_{i=1}^k A_i^c = \left(\bigcup_{i=1}^k A_i\right)^c$$
Without the outer bounding canvas, $\bigcup_{\vec{b} \neq \vec{0}} M_{\vec{b}} \neq U$ whenever there exist elements outside the union, violating the definition of a partition. In Monster Sort, any creature with $\vec{b}(x) = \vec{0}$ must be placed in the outer canvas, reinforcing the **Principle of the Universal Complement**.

\subsection{De Morgan's Laws and Planar Topological Limits}
\begin{itemize}
  \item \textbf{De Morgan's Duality:} In World 4 (Mirror Peaks), players solve expressions like $(A \cup B)^c$. The game proves $(A \cup B)^c = A^c \cap B^c$ by verifying that dragging creatures into the region satisfying $\neg(P_A \lor P_B)$ selects identical bitmasks as $(1 - P_A) \land (1 - P_B)$.
  \item \textbf{Topological Limit (Grünbaum / Venn Theorem):} In the plane, 3 circles can form all $2^3 = 8$ connected regions. However, for $k = 4$, circular curves can produce at most 14 regions ($r \le n^2 - n + 2$), which is strictly less than $2^4 = 16$. This explains why World 3 caps standard circular Venn diagrams at $k=3$ before introducing higher-order compound predicate chambers.
\end{itemize}

\newpage
\section{Mathematical Reasoning: Realm 2 --- The Robot Maze Challenge}

\subsection{Graph-Theoretic Formalization of Mazes}
A visual maze is mathematically formalized as an **unweighted undirected planar graph** $G = (V, E)$:
\begin{definition}[Grid Graph Representation]
Given a 2D lattice of dimensions $R \times C$:
\begin{itemize}
  \item \textbf{Vertex Set $V$:} Every open (walkable) cell $(r, c)$ constitutes a vertex:
  $$V = \{(r, c) \in \{0, \dots, R-1\} \times \{0, \dots, C-1\} \mid \text{maze}[r][c] \neq \text{`\#'}\}$$
  \item \textbf{Edge Set $E$:} An edge joins two vertices if and only if they are orthogonal unit neighbors (Manhattan distance 1):
  $$E = \{(u, v) \in V \times V \mid u = (r, c), v = (r', c'), \ |r - r'| + |c - c'| = 1\}$$
\end{itemize}
\end{definition}

\begin{lemma}[The Handshake Lemma in Grid Graphs]
For every maze level:
$$\sum_{v \in V} \operatorname{deg}(v) = 2|E|$$
Each interior cell has degree $\operatorname{deg}(v) \le 4$, edge cells $\le 3$, and corner cells $\le 2$. Walls act as induced subgraph vertex deletions ($G' = G[V \setminus V_{\text{obs}}]$).
\end{lemma}

\subsection{State-Space Representation and Geodesic Shortest Path}
\begin{definition}[Path and Geodesic Distance]
A path $\pi$ of length $L$ from start vertex $s = (r_0, c_0)$ to destination $d = (r_g, c_g)$ is a sequence of vertices:
$$\pi = \langle v_0, v_1, \dots, v_L \rangle \quad \text{such that } v_0 = s, \ v_L = d, \ \text{and } (v_i, v_{i+1}) \in E \ \forall i$$
The **geodesic distance** $\delta(s, d)$ is defined as:
$$\delta(s, d) = \begin{cases} \min \{L \mid \text{there exists a path } \pi \text{ of length } L \text{ from } s \text{ to } d\} \\ \infty, \quad \text{if no path exists} \end{cases}$$
\end{definition}

\subsection{Proof of Breadth-First Search (BFS) Optimality}
\begin{theorem}[BFS Shortest-Path Correctness on Unweighted Graphs]
On a finite graph $G = (V, E)$ with uniform edge weights $w(e) = 1$, Breadth-First Search initialized at $s$ computes $d[v] = \delta(s, v)$ for all $v \in V$ in $O(|V| + |E|)$ time.
\end{theorem}

\begin{proof}[Layer Induction Proof (Examiner-Ready)]
Define distance layers $L_d = \{v \in V \mid \delta(s, v) = d\}$ for $d \ge 0$, where $L_0 = \{s\}$.

\textbf{Induction Hypothesis:} BFS dequeues every vertex in layer $L_d$ before dequeuing any vertex in layer $L_{d+1}$, and records distance $d[v] = d = \delta(s, v)$ for every $v \in L_d$.

\textbf{Base Case ($d=0$):} At start, $Q = \langle s \rangle$ with $d[s] = 0 = \delta(s, s)$. $s$ is dequeued before any other vertex. The base case holds.

\textbf{Inductive Step:} Assume true for layers $0, 1, \dots, d$.
Consider any vertex $w \in L_{d+1}$. By definition, $w$ has at least one neighbor $v \in L_d$, and $w \notin L_{\le d}$.
Because $w \notin L_{\le d}$, $w$ has not been discovered prior to processing layer $L_d$.
When BFS dequeues neighbor $v \in L_d$ (which occurs during layer $d$ processing by induction):
\begin{enumerate}
  \item $w$ is undiscovered, so BFS assigns $d[w] = d[v] + 1 = d + 1 = \delta(s, w)$.
  \item BFS sets $\operatorname{parent}[w] = v$ and enqueues $w$ at the tail of the FIFO queue $Q$.
\end{enumerate}
Because $Q$ is First-In-First-Out, all vertices of $L_{d+1}$ are enqueued behind all vertices of $L_d$, and before any vertex discovered from $L_{d+1}$.
Thus, all vertices in $L_{d+1}$ are dequeued strictly after $L_d$ and strictly before $L_{d+2}$.
Tracing the parent chain back from $d$ reconstructs an optimal path $\pi^*$ of length exactly $d[d] = \delta(s, d)$.
\end{proof}

\subsection{Why Plain BFS Fails on Weighted Graphs (The Mud Counterexample)}
In the interactive lesson (Learn F), we demonstrate what happens if terrain has unequal costs (e.g., mud tiles cost 3 steps, paved tiles cost 1 step).
\begin{itemize}
  \item BFS measures **edge count**, not **edge cost**.
  \item A path with 3 paved steps has total cost $1 + 1 + 1 = 3$.
  \item A path with 2 muddy steps has total cost $3 + 3 = 6$, yet BFS would mistakenly choose it because $2 < 3$ in hop count!
  \item This directly teaches students why BFS requires unit-cost edges, motivating **Dijkstra's Algorithm** for weighted graphs.
\end{itemize}

\subsection{Multiple Shortest Paths and Dynamic Programming}
In Level 11 (``Many Shortest Paths''), there are multiple distinct paths of minimal length $\delta(s, d) = 14$. The game counts them using Dynamic Programming over the BFS Directed Acyclic Graph (DAG):
$$\text{numPaths}(v) = \sum_{u \in \operatorname{Adj}(v) \cap L_{d(v)-1}} \text{numPaths}(u) \quad \text{with } \text{numPaths}(s) = 1$$
This yields exactly 16 shortest paths, illustrating the combinatorial nature of grid DAGs.

\newpage
\section{The Unifying Mathematical Duality}

The faculty examiner will ask: \textit{``How do these two projects relate mathematically? Why combine them?''}

\subsection{Static Classification vs.\ Dynamic State-Space Search}
\begin{table}[htbp]
\centering
\small
\begin{tabularx}{\textwidth}{@{}P{3.5cm}Y Y@{}}
\toprule
\textbf{Mathematical Dimension} & \textbf{Monster Sort (Set Theory Realm)} & \textbf{Robot Maze (Graph Theory Realm)} \\
\midrule
\textbf{Primitive Objects} & Monsters with feature vectors $\vec{f} \in \mathcal{F}$ & Grid cells as vertices $V \subset \mathbb{Z}^2$ \\
\textbf{Fundamental Structure} & Predicates $P_i : U \to \{0, 1\}$ and sets $A_i \subseteq U$ & Orthogonal adjacency relation $E \subseteq V \times V$ \\
\textbf{Core Problem} & \textbf{Partitioning Problem:} Map each element to its unique minterm region in $O(k)$ & \textbf{Optimization Problem:} Find minimum edge sequence $\pi^* = \langle s, \dots, d \rangle$ in $O(|V|+|E|)$ \\
\textbf{Decision Tool} & Characteristic Bitmask Vector: $\vec{b}(x) \in \{0, 1\}^k$ & FIFO Queue Frontier Expansion: $Q = \langle v_1, \dots, v_m \rangle$ \\
\textbf{Failure Mode} & Set Constraint Violation (placing $x \notin A$ inside circle $A$) & Obstacle Collision (stepping into $V_{\text{obs}}$) or sub-optimal cycles \\
\bottomrule
\end{tabularx}
\end{table}

\subsection{The Deep Duality: Graph Traversal is Dynamic Set Partitioning}
At any instantaneous step $t$ of the Breadth-First Search algorithm in Robot Maze, the graph's vertex set $V$ is partitioned into three pairwise disjoint subsets:
$$V = \operatorname{Visited}^{(t)} \cup \operatorname{Frontier}^{(t)} \cup \operatorname{Unvisited}^{(t)}$$
\begin{enumerate}
  \item $\operatorname{Visited}^{(t)}$: Vertices whose shortest distances are permanently resolved.
  \item $\operatorname{Frontier}^{(t)}$ (the FIFO Queue): Vertices discovered at current wavefront $L_d$ or $L_{d+1}$.
  \item $\operatorname{Unvisited}^{(t)}$: Vertices whose state is yet to be discovered.
\end{enumerate}
Every step of BFS moves a vertex from $\operatorname{Frontier}$ to $\operatorname{Visited}$ and its undiscovered neighbors from $\operatorname{Unvisited}$ to $\operatorname{Frontier}$.
\textbf{Thus, state-space exploration is mathematically equivalent to iteratively evolving a 3-block set partition of $V$.}

\subsection{The Sorting Hat Cognitive Profiling Algorithm}
The Sorting Hat questionnaire maps 5 situational questions into cognitive vectors $\vec{s} = (s_{\text{set}}, s_{\text{graph}})$.
\begin{itemize}
  \item \textbf{Taxonomic / Categorical Affinity:} Tendency to solve problems by clustering, identifying invariant traits, and filtering (Boolean predicates $\to$ Monster Sort).
  \item \textbf{Topological / Procedural Affinity:} Tendency to solve problems by chaining sequences, exploring state transitions, and optimizing paths (Graph relations $\to$ Robot Maze).
  \item \textbf{Dual Mastery:} Balanced scores unlock both realms immediately, demonstrating that computer science mastery requires fluency in both categorical data classification and procedural algorithmic navigation.
\end{itemize}

\newpage
\section{Anticipated Viva-Voce Questions \& Model Answers}

This section contains the exact questions the professor (``mam'') is most likely to ask during the project viva, paired with comprehensive, mathematically sound model answers.

\begin{vivabox}
\textbf{Q1: Why does a Venn diagram with $k$ sets partition space into exactly $2^k$ regions?}

\textbf{Model Answer:}
``Each set $A_i$ is defined by a unary Boolean predicate $P_i$. By the Law of Excluded Middle, any element $x \in U$ either belongs to $A_i$ ($P_i(x) = 1$) or does not ($P_i(x) = 0$). For $k$ independent predicates, each element has a characteristic bit-vector $\vec{b}(x) = (P_1(x), \dots, P_k(x)) \in \{0, 1\}^k$. By the Rule of Product in combinatorics, the number of distinct binary sequences of length $k$ is $2 \times 2 \times \dots \times 2 = 2^k$. Each bit sequence corresponds to one fundamental minterm intersection $M_{\vec{b}} = \bigcap_{i=1}^k A_i^{(b_i)}$, so there are exactly $2^k$ possible mutually exclusive regions.''
\end{vivabox}

\begin{vivabox}
\textbf{Q2: Why can't we draw a standard Venn diagram with 4 circles in the plane?}

\textbf{Model Answer:}
``This is governed by Venn's Theorem and Branko Gr{\"u}nbaum's geometric investigations. In the Euclidean plane, any two circles intersect in at most 2 points. By planar graph theory and Euler's formula $V - E + F = 2$, $n$ circles in general position can divide the plane into at most $n^2 - n + 2$ connected regions. For $n = 3$, $3^2 - 3 + 2 = 8 = 2^3$, which works. But for $n = 4$, $4^2 - 4 + 2 = 14 < 2^4 = 16$. Two required minterms cannot be geometrically formed using circles. To represent $k=4$, one must use non-convex shapes, ellipses (as John Venn did), or higher-dimensional hypercubes.''
\end{vivabox}

\begin{vivabox}
\textbf{Q3: How does Monster Sort demonstrate De Morgan's Laws?}

\textbf{Model Answer:}
``In World 4 (Mirror Peaks), the game presents expressions such as $(A \cup B)^c$.
Mathematically, De Morgan's Law states:
$$(A \cup B)^c = A^c \cap B^c$$
In the game, a monster is placed into the outer universal realm if it satisfies $\neg(P_A \lor P_B)$. By Boolean truth tables, $\neg(P_A \lor P_B) \equiv \neg P_A \land \neg P_B$. The game engine computes the indicator function $\chi_{(A \cup B)^c}(x) = 1 - \max(\chi_A(x), \chi_B(x))$ and confirms it evaluates to 1 if and only if $(1 - \chi_A(x)) \cdot (1 - \chi_B(x)) = 1$. The visual interface highlights both representations as identical target chambers.''
\end{vivabox}

\begin{vivabox}
\textbf{Q4: In Robot Maze, why does Breadth-First Search guarantee the shortest path while Depth-First Search (DFS) does not?}

\textbf{Model Answer:}
``BFS explores vertices in strictly increasing order of distance using a FIFO queue. It processes all vertices at distance $d$ (layer $L_d$) before visiting any vertex at distance $d+1$. Therefore, the very first time the destination vertex $d$ is discovered, the path leading to it is guaranteed to have minimal edge count $\delta(s, d)$. In contrast, DFS uses a LIFO stack (or recursion) and explores deep down a single branch. DFS will follow arbitrary corridors and can reach the goal via a highly convoluted, circuitous path of length $L \gg \delta(s, d)$ before backtracking, failing to guarantee shortest path length.''
\end{vivabox}

\begin{vivabox}
\textbf{Q5: Why does BFS fail to find the shortest path if grid tiles have different traversal costs?}

\textbf{Model Answer:}
``The correctness proof of BFS relies fundamentally on the **Unit Cost Axiom**: every edge has weight $w(e) = 1$. This ensures that number of edge hops equals total path cost. If a grid contains mud tiles with weight $w = 3$ and road tiles with weight $w = 1$, a 2-hop mud path has cost $3 + 3 = 6$, while a 3-hop road path has cost $1 + 1 + 1 = 3$. Standard BFS only counts hops, so it would dequeue the 2-hop mud path first, which is suboptimal. For non-uniform non-negative edge costs, one must generalize BFS to **Dijkstra's Algorithm** using a priority queue indexed by cumulative distance.''
\end{vivabox}

\begin{vivabox}
\textbf{Q6: What is the time and space complexity of both games' core algorithms?}

\textbf{Model Answer:}
``\textbf{1. Monster Sort Placement:} For $k$ predicates and $|U|$ monsters, extracting features and evaluating the $k$-bit characteristic vector takes $O(k)$ time per monster. Direct table lookup to the target minterm is $O(1)$. Total verification time for a level is $O(k \cdot |U|)$. Space complexity is $O(2^k + |U|)$ to store the $2^k$ regions and monster attributes.
\textbf{2. Robot Maze BFS:} On an $R \times C$ maze graph $G=(V, E)$, $|V| \le RC$ and $|E| < 2RC$. Each vertex is enqueued and dequeued at most once ($O(|V|)$ operations). For each dequeued vertex, its orthogonal neighbors (at most 4) are inspected ($O(|E|)$ operations). Thus, time complexity is $O(|V| + |E|) = O(RC)$. Space complexity is $O(|V|)$ to store the visited boolean array, parent pointers, and FIFO queue.''
\end{vivabox}

\begin{vivabox}
\textbf{Q7: How did you implement the queue in Robot Maze to ensure $O(1)$ operations?}

\textbf{Model Answer:}
``In JavaScript, using `Array.prototype.shift()` to dequeue from the front takes $O(n)$ time because it re-indexes every remaining element in memory, degrading total BFS time from $O(V)$ to $O(V^2)$. To prevent this, in `js/bfs.js`, we implemented the queue as an array with an advancing `head` index pointer. Dequeueing simply reads `queue[head++]`, which is strictly $O(1)$ in time. This guarantees that our browser implementation maintains the theoretical $O(|V| + |E|)$ asymptotic complexity.''
\end{vivabox}

\begin{vivabox}
\textbf{Q8: How does the game verify the Handshake Lemma?}

\textbf{Model Answer:}
``By Euler's Handshake Lemma, for any undirected graph, $\sum_{v \in V} \operatorname{deg}(v) = 2|E|$. In our automated test suite (`tests/run-tests.js`), we iterate over every open cell in each of the 12 levels, compute its degree using orthogonal neighbor checks (`getNeighbors()`), sum the degrees across all vertices, and independently count the bidirectional edges in `listEdges()`. The test programmatically asserts that the sum of degrees is exactly equal to twice the total edge count for all levels, proving graph structural consistency.''
\end{vivabox}

\begin{vivabox}
\textbf{Q9: In Monster Sort, what is the role of the benchmark dataset from Section 10?}

\textbf{Model Answer:}
``Section 10 defines a canonical ground-truth test suite on universe $U = \{M_1, \dots, M_6\}$ with predicates Red ($R$), Horns ($H$), and Wings ($W$). It rigorously validates all Boolean set operators:
\begin{itemize}
  \item Intersection: $R \cap H = \{M_1, M_6\}$
  \item Union: $R \cup H = \{M_1, M_2, M_3, M_6\}$
  \item Complement: $U \setminus R = \{M_2, M_4, M_5\}$
  \item Set Difference: $R \setminus H = \{M_3\}$
  \item Triple Intersection: $R \cap H \cap W = \{M_6\}$
  \item Compound Difference: $(R \cup H) \setminus W = \{M_1\}$
  \item Universal Complement: $(R \cup H \cup W)^c = \{M_4\}$
\end{itemize}
Our Vitest test suite runs automated assertions verifying that our game engine's set operations match these exact outputs with zero discrepancies.''
\end{vivabox}

\begin{vivabox}
\textbf{Q10: What is a Cut Vertex in Robot Maze, and how does it relate to graph connectivity?}

\textbf{Model Answer:}
``In graph theory, a cut vertex (or articulation point) is a vertex whose removal increases the number of connected components of graph $G$. In the Robot Maze 'Math Detective' lab (Experiment 3), students place wall obstacles on single-width corridors. If that corridor cell is a cut vertex separating start $s$ and destination $d$, its removal disconnects $G$ into two disjoint components $G_1$ and $G_2$ with $s \in G_1$ and $d \in G_2$. BFS exploration terminates with an empty queue, returning $\delta(s, d) = \infty$. This visually demonstrates vertex connectivity and network reliability.''
\end{vivabox}

\begin{vivabox}
\textbf{Q11: How does the Principle of Inclusion-Exclusion (PIE) relate to Monster Sort?}

\textbf{Model Answer:}
``Because the $2^k$ minterms form a pairwise disjoint partition of the universe, the size of any compound set $S$ is strictly the sum of the cardinalities of its constituent minterms: $|S| = \sum_{\vec{b} \in \mathcal{I}_S} |M_{\vec{b}}|$. For the union of three sets $R, H, W$:
$$|R \cup H \cup W| = |R| + |H| + |W| - (|R \cap H| + |R \cap W| + |H \cap W|) + |R \cap H \cap W|$$
This algebraic formula arises directly from the fact that overlapping regions are counted multiple times in the sum $|R| + |H| + |W|$ and must be compensated by alternating subtractions and additions. In Monster Sort, counting monsters inside the Venn rings visually demonstrates why PIE is true.''
\end{vivabox}

\begin{vivabox}
\textbf{Q12: What were the individual technical and mathematical contributions of each member?}

\textbf{Model Answer:}
\begin{itemize}
  \item \textbf{Dungar Soni (B23CS1105):} Lead Full-Stack Architect; designed the Next.js 16/React 19 game engine, state machines, drag-and-drop mechanics, Vercel deployment, and GitHub pipeline.
  \item \textbf{Radtiya Saraf (B23CS1107):} Theoretical Lead; authored formal proofs for the Venn Partition Theorem and BFS Geodesic Optimality, and formalized characteristic vector mathematics.
  \item \textbf{Prakhar Goyal (B23CS1106):} Geometric Modeling; computed planar SVG clipping masks and coordinates for the 8 Venn regions, and modeled grid graph adjacency and wall collision physics.
  \item \textbf{Arafat (B23CS1103):} Asset & Audio Engineering; created the 7-attribute monster schema, vector SVG artwork, procedural Web Audio tone synthesis, and accessible CSS styling.
  \item \textbf{Bhawani (B23CS1104):} Pedagogical Logic & Experiments; created the Sorting Hat personality quiz scoring algorithm, Robot Race gameplay, and designed the school outreach lesson plan.
\end{itemize}
\end{vivabox}

\end{document}
